\section{Step 37 Grade the public behavior independently of the illustration}\label{pcm-s37}

The construction observations name temporary registers because they explain this particular sequence. Complete behavioral grading instead checks the callable interface, exact results, preserved state and every caller-data access. It must reject an out-of-slice store even if a later store restores the byte. In free writing, choose different temporaries, helper labels or a scalar implementation while retaining the public contract.

\subsection{Expected observation}
The completed scalar and MSA implementations pass the same behavioral cases. The untouched starter is rejected. Exact instruction counts are measurements of the examples, not grading requirements.

\Needspace{11\baselineskip}
\subsection{Check your understanding}
Can a correct alternative use different temporary registers?

\begin{itemize}
  \item \textbf{A} No, because the observation table names them.
  \item \textbf{B} Yes, if the public contract still holds.
  \item \textbf{C} Only if it also copies the example’s instruction count.
\end{itemize}
